% GENERATED FILE. Edit canonical lessons, then run npm run generate:books.
% canonical-source-hash: fnv1a64:f0fe03f6112c7e9d
% canonical-lessons: JA-C43-suna, JA-C43-tanbo, JA-C43-hoshi, JA-C43-kumo, JA-C43-taiyou

\chapter{Sand, Rice field, Star, Cloud, Sun}
\label{ch:ja-sand-rice-field-star-cloud-sun}
\hlchaptermodality{43}

\begin{chapteropening}
\textbf{By the end of this chapter:} \emph{I can say sand, a rice field, a star, a cloud and the sun.}

Complete the last lesson of chapter 43: i can say sand, a rice field, a star, a cloud and the sun.
\end{chapteropening}

\section[suna]{\ja{すな} (suna) — sand}

\label{lesson:JA-C43-suna}



\begin{quote}
Before the new one: say the Japanese for a forest, then the Japanese for a stone.
\end{quote}

\subsection*{You'll want to know: \ja{すな}}
\textbf{\ja{すな}} — \emph{suna} — \textquotedblleft{}sand\textquotedblright{}.

\begin{grammarlens}[title={What you've built}]
One more word for this chapter.
\end{grammarlens}

\subsection*{Your turn}
\begin{itemize}
\raggedright
  \item \emph{Say it:} \emph{suna}
  \item \emph{Say it:} \emph{suna}, once more
  \item \emph{Say it:} say \emph{suna}
  \item \emph{Recall:} say the Japanese for a forest, then the Japanese for a stone, then say \emph{suna} again
\end{itemize}

\begin{tcolorbox}[breakable,colback=teal!4,colframe=teal!35!black,title={Before you move on}]
What does \ja{すな} mean? (\textquotedblleft{}Sand\textquotedblright{}.) Say it once more.
\end{tcolorbox}
\section[tanbo]{\ja{たんぼ} (tanbo) — a rice field}

\label{lesson:JA-C43-tanbo}



\begin{quote}
Before the new one: say the Japanese for a stone, then the Japanese for sand.
\end{quote}

\subsection*{You'll want to know: \ja{たんぼ}}
\textbf{\ja{たんぼ}} — \emph{tanbo} — \textquotedblleft{}a rice field\textquotedblright{}.

\begin{grammarlens}[title={What you've built}]
One more word for this chapter.
\end{grammarlens}

\subsection*{Your turn}
\begin{itemize}
\raggedright
  \item \emph{Say it:} \emph{tanbo}
  \item \emph{Say it:} \emph{tanbo}, once more
  \item \emph{Say it:} say \emph{tanbo}
  \item \emph{Recall:} say the Japanese for a stone, then the Japanese for sand, then say \emph{tanbo} again
\end{itemize}

\begin{tcolorbox}[breakable,colback=teal!4,colframe=teal!35!black,title={Before you move on}]
What does \ja{たんぼ} mean? (\textquotedblleft{}A rice field\textquotedblright{}.) Say it once more.
\end{tcolorbox}
\section[hoshi]{\ja{ほし} (hoshi) — a star}

\label{lesson:JA-C43-hoshi}



\begin{quote}
Before the new one: say the Japanese for sand, then the Japanese for a rice field.
\end{quote}

\subsection*{You'll want to know: \ja{ほし}}
\textbf{\ja{ほし}} — \emph{hoshi} — \textquotedblleft{}a star\textquotedblright{}.

\begin{grammarlens}[title={What you've built}]
One more word for this chapter.
\end{grammarlens}

\subsection*{Your turn}
\begin{itemize}
\raggedright
  \item \emph{Say it:} \emph{hoshi}
  \item \emph{Say it:} \emph{hoshi}, once more
  \item \emph{Say it:} say \emph{hoshi}
  \item \emph{Recall:} say the Japanese for sand, then the Japanese for a rice field, then say \emph{hoshi} again
\end{itemize}

\begin{tcolorbox}[breakable,colback=teal!4,colframe=teal!35!black,title={Before you move on}]
What does \ja{ほし} mean? (\textquotedblleft{}A star\textquotedblright{}.) Say it once more.
\end{tcolorbox}
\section[kumo]{\ja{くも} (kumo) — a cloud; a spider}

\label{lesson:JA-C43-kumo}



\begin{quote}
Before the new one: say the Japanese for a rice field, then the Japanese for a star.
\end{quote}

\subsection*{You'll want to know: \ja{くも}}
\textbf{\ja{くも}} — \emph{kumo} — \textquotedblleft{}a cloud; a spider\textquotedblright{}.

\begin{grammarlens}[title={What you've built}]
One more word for this chapter.
\end{grammarlens}

\subsection*{Your turn}
\begin{itemize}
\raggedright
  \item \emph{Say it:} \emph{kumo}
  \item \emph{Say it:} \emph{kumo}, once more
  \item \emph{Say it:} say \emph{kumo}
  \item \emph{Recall:} say the Japanese for a rice field, then the Japanese for a star, then say \emph{kumo} again
\end{itemize}

\begin{tcolorbox}[breakable,colback=teal!4,colframe=teal!35!black,title={Before you move on}]
What does \ja{くも} mean? (\textquotedblleft{}A cloud; a spider\textquotedblright{}.) Say it once more.
\end{tcolorbox}
\section[taiyou]{\ja{たいよう} (taiyou) — the sun}

\label{lesson:JA-C43-taiyou}



\begin{quote}
Before the new one: say the Japanese for a star, then the Japanese for a cloud; a spider.
\end{quote}

\subsection*{You'll want to know: \ja{たいよう}}
\textbf{\ja{たいよう}} — \emph{taiyou} — \textquotedblleft{}the sun\textquotedblright{}.

\begin{grammarlens}[title={What you've built}]
One more word for this chapter.
\end{grammarlens}

\subsection*{Your turn}
\begin{itemize}
\raggedright
  \item \emph{Say it:} \emph{taiyou}
  \item \emph{Say it:} \emph{taiyou}, once more
  \item \emph{Say it:} say \emph{taiyou}
  \item \emph{Recall:} say the Japanese for a star, then the Japanese for a cloud; a spider, then say \emph{taiyou} again
\end{itemize}

\begin{tcolorbox}[breakable,colback=teal!4,colframe=teal!35!black,title={Before you move on}]
What does \ja{たいよう} mean? (\textquotedblleft{}The sun\textquotedblright{}.) Say it once more.
\end{tcolorbox}
